\chapter{Contexto do Projeto}
\label{cap:contexto}

\orientacao{Este capítulo apresenta o contexto real no qual o projeto está inserido, caracterizando o cenário, o público envolvido, as necessidades identificadas e as regras de negócio do domínio. Sua função é estabelecer a compreensão do problema sob a perspectiva da comunidade atendida antes de apresentar a solução computacional.}

\orientacao{Incluir estudo sobre a área de atuação do parceiro do projeto, mercado, índices, demandas, etc.}

\section{Cenário}
\orientacao{Descrever o contexto em que o projeto será aplicado, incluindo as características relevantes do ambiente, organização, atividade ou situação que originou a demanda.}

\section{Público-Alvo}
\orientacao{Caracterizar as pessoas, grupos ou organizações que utilizarão, serão atendidos ou serão diretamente afetados pelo projeto.}

\section{Necessidades Identificadas}
\orientacao{Apresentar as necessidades observadas no contexto do projeto que motivaram a proposição da solução, sem ainda detalhar os requisitos do sistema.}

\section{Regras de Negócio}
\orientacao{Enunciar e discutir as regras, condições, políticas e restrições próprias do domínio que devem ser respeitadas pela solução, mantendo a descrição em nível conceitual e independente da implementação.}

\begin{landscape}
\begin{figure}[htbp]
    \centering
    \caption{Diagrama de atividades do cadastro e da publicação de anúncio}
    \label{fig:atividades}
    \resizebox{\linewidth}{!}{\begin{tikzpicture}[
  font=\footnotesize,
  acao/.style={draw, rounded corners=4pt, text width=2.3cm, align=center, minimum height=1cm, inner sep=3pt,
    execute at begin node={\hyphenpenalty=10000}},
  ponto/.style={draw, diamond, minimum size=0.55cm, inner sep=0pt},
  inicio/.style={circle, fill=black, minimum size=0.35cm, inner sep=0pt},
  final/.style={circle, draw, minimum size=0.4cm, inner sep=0pt},
  fluxo/.style={-{Stealth[length=2mm]}},
  guarda/.style={font=\scriptsize},
]
  % raias
  \draw (0,0) rectangle (23.8,9.6);
  \draw (0,6.0) -- (23.8,6.0);
  \draw (0,3.2) -- (23.8,3.2);
  \draw (0.8,0) -- (0.8,9.6);
  \node[rotate=90] at (0.4,7.8) {Vendedor};
  \node[rotate=90] at (0.4,4.6) {\sistema};
  \node[rotate=90] at (0.4,1.6) {Marketplace};

  % cadastro
  \node[inicio] (ini) at (1.4,8.4) {};
  \node[acao] (acessar) at (3.2,8.4) {Acessar o \sistema};
  \node[acao] (produto) at (6.3,8.4) {Cadastrar produto};
  \node[acao] (anuncio) at (9.4,8.4) {Cadastrar anúncio selecionando o marketplace};

  % verificacao de anuncio existente
  \node[acao] (verificar) at (9.4,4.7) {Verificar anúncio existente no marketplace};
  \node[ponto] (d0) at (11.5,4.7) {};
  \node[guarda, anchor=north, align=center] at (d0.south) {anúncio já\\existe?};
  \node[acao] (informar) at (13.8,4.7) {Informar anúncio já existente};
  \node[final] (fim1) at (15.8,4.7) {};
  \fill (fim1.center) circle (0.12cm);

  % envio, validacao e publicacao
  \node[ponto] (m1) at (16.8,8.4) {};
  \node[acao] (enviar) at (19.0,8.4) {Enviar anúncio};
  \node[acao] (corrigir) at (16.8,6.9) {Corrigir anúncio};
  \node[acao] (validar) at (19.0,2.3) {Validar anúncio};
  \node[ponto] (d1) at (16.8,2.3) {};
  \node[guarda, anchor=east] at (d1.west) {anúncio recusado?};
  \node[acao] (publicar) at (21.2,0.8) {Publicar anúncio};
  \node[final] (fim2) at (23.2,0.8) {};
  \fill (fim2.center) circle (0.12cm);

  \draw[fluxo] (ini) -- (acessar);
  \draw[fluxo] (acessar) -- (produto);
  \draw[fluxo] (produto) -- (anuncio);
  \draw[fluxo] (anuncio) -- (verificar);
  \draw[fluxo] (verificar) -- (d0);
  \draw[fluxo] (d0) -- node[guarda, above] {[sim]} (informar);
  \draw[fluxo] (informar) -- (fim1);
  \draw[fluxo] (d0) |- node[guarda, pos=0.25, right] {[não]} (m1);
  \draw[fluxo] (m1) -- (enviar);
  \draw[fluxo] (enviar) -- (validar);
  \draw[fluxo] (validar) -- (d1);
  \draw[fluxo] (d1) -- node[guarda, right] {[sim]} (corrigir);
  \draw[fluxo] (corrigir) -- (m1);
  \draw[fluxo] (d1) |- node[guarda, pos=0.25, left] {[não]} (publicar);
  \draw[fluxo] (publicar) -- (fim2);
\end{tikzpicture}
}
    \fonte{Elaborado pelos autores.}
\end{figure}
\end{landscape}
